\documentclass[10pt]{article}
\providecommand{\FIGDIR}{fig}
\newcommand{\DIGESTDATE}{30 August 2026}
% ===================== PM Research Watch - shared preamble =====================
\usepackage[T1]{fontenc}
\usepackage[utf8]{inputenc}
\usepackage{amsmath,amssymb}
\usepackage{tgpagella}
\usepackage{tgheros}
\renewcommand{\sfdefault}{qhv}
\renewcommand{\ttdefault}{lmtt}
\usepackage{microtype}
\usepackage[english]{babel}

\usepackage[a4paper,top=17mm,bottom=17mm,left=17mm,right=17mm,
            headsep=6mm,footskip=10mm]{geometry}
\usepackage{graphicx}
\usepackage{xcolor}
\usepackage{colortbl}
\usepackage{tikz}
\usetikzlibrary{calc}
\usepackage[most]{tcolorbox}
\usepackage{titlesec}
\usepackage{fancyhdr}
\usepackage{booktabs}
\usepackage{array}
\usepackage{longtable}
\usepackage{enumitem}
\usepackage{multicol}
\usepackage{ragged2e}
\usepackage[hidelinks]{hyperref}

% ------------------------------------------------------------------ palette
\definecolor{Ink}    {HTML}{14262E}
\definecolor{Deep}   {HTML}{0F4C5C}
\definecolor{Teal}   {HTML}{1F7A8C}
\definecolor{Sky}    {HTML}{4FA3C4}
\definecolor{Sage}   {HTML}{6FA287}
\definecolor{Amber}  {HTML}{E3A008}
\definecolor{Clay}   {HTML}{B5643C}
\definecolor{Coral}  {HTML}{D1495B}
\definecolor{Violet} {HTML}{6A5B9E}
\definecolor{Slate}  {HTML}{4A6572}
\definecolor{Paper}  {HTML}{FBF9F5}
\definecolor{Mist}   {HTML}{EEEAE2}
\definecolor{Rule}   {HTML}{D9D3C9}

\pagecolor{Paper}
\color{Ink}

\hypersetup{colorlinks=true, linkcolor=Deep, urlcolor=Teal, citecolor=Teal,
            pdftitle={PM Research Watch}, pdfauthor={Automated daily digest}}

% ------------------------------------------------------------------ headings
\titleformat{\section}
  {\sffamily\bfseries\large\color{Deep}}{}{0pt}
  {}[{\vspace{1.2mm}\color{Amber}\titlerule[1.1pt]}]
\titlespacing*{\section}{0pt}{5.5mm}{2.8mm}

\titleformat{\subsection}
  {\sffamily\bfseries\normalsize\color{Teal}}{}{0pt}{}
\titlespacing*{\subsection}{0pt}{3.6mm}{1.4mm}

\setlist[itemize]{leftmargin=4.2mm, itemsep=1.1mm, topsep=1.3mm, parsep=0pt,
                  label=\textcolor{Teal}{\footnotesize$\blacktriangleright$}}

% ------------------------------------------------------------------ header/footer
\pagestyle{fancy}
\fancyhf{}
\renewcommand{\headrulewidth}{0pt}
\renewcommand{\footrulewidth}{0pt}
\fancyhead[L]{\sffamily\scriptsize\color{Slate}PM RESEARCH WATCH\;\textcolor{Rule}{\textbar}\;\DIGESTDATE}
\fancyhead[R]{\sffamily\scriptsize\color{Slate}Particulate matter monitoring \& health effects}
\fancyfoot[C]{\sffamily\scriptsize\color{Slate}\thepage}

% ------------------------------------------------------------------ components
% metric tile
\newcommand{\metric}[3]{%
  \begin{tcolorbox}[enhanced, colback=white, colframe=Rule, boxrule=0.5pt,
    arc=1.6mm, left=2.2mm, right=2.2mm, top=1.8mm, bottom=1.8mm,
    borderline west={1.7mm}{0pt}{#3}, nobeforeafter, valign=center,
    height=17mm, width=\linewidth]
    \RaggedRight
    {\sffamily\bfseries\fontsize{17}{17}\selectfont\color{#3}#1}\\[0.5mm]
    {\sffamily\fontsize{6.5}{7.4}\selectfont\color{Slate}#2}
  \end{tcolorbox}}

% pull-quote / key insight box
\newtcolorbox{keybox}[1][]{enhanced, colback=Mist, colframe=Mist,
  arc=1.6mm, left=3.2mm, right=3.2mm, top=2.4mm, bottom=2.4mm,
  borderline west={1.7mm}{0pt}{Amber}, #1}

% subtopic band
\newcommand{\band}[2]{%
  \vspace{3.2mm}\noindent
  \begin{tcolorbox}[enhanced, colback=#2, colframe=#2, arc=1.2mm,
    left=2.6mm, right=2.6mm, top=1.1mm, bottom=1.1mm, boxrule=0pt,
    width=\linewidth, nobeforeafter]
    {\sffamily\bfseries\small\color{white}#1}
  \end{tcolorbox}\vspace{1.6mm}}

% one paper entry: {short cite}{title}{journal}{doi}{body}
\newcommand{\paperentry}[5]{%
  \par\vspace{1.9mm}\noindent
  {\sffamily\bfseries\footnotesize\color{Ink}#1}\;%
  {\sffamily\fontsize{7}{8}\selectfont\color{Slate}#3}\\[0.4mm]
  {\footnotesize\itshape\color{Deep}#2}\\[0.6mm]
  {\footnotesize #4}\\[0.3mm]
  {\fontsize{6.6}{7.6}\selectfont\color{Slate}\texttt{doi:}\,\href{https://doi.org/#5}{\texttt{\detokenize{#5}}}}
  \par}

% caption style
\newcommand{\figcap}[1]{\par\vspace{1.0mm}
  {\fontsize{7.2}{8.6}\selectfont\color{Slate}\textsf{\textbf{Fig.}}\; #1\par}\vspace{1.5mm}}

\providecommand{\permil}{\textperthousand}

\setlength{\parindent}{0pt}
\setlength{\parskip}{0pt}
\renewcommand{\arraystretch}{1.18}


\begin{document}

% ==================================================================== MASTHEAD
\thispagestyle{empty}
\begin{tikzpicture}[remember picture, overlay]
  \fill[Deep] (current page.north west) rectangle
      ($(current page.north east)+(0,-67mm)$);
  \fill[Amber] ($(current page.north west)+(0,-67mm)$) rectangle
      ($(current page.north east)+(0,-68.4mm)$);
  \foreach \x/\y/\r in {22/12/0.9, 41/26/1.5, 63/9/0.7, 88/31/1.1, 112/17/1.9,
                        138/29/0.8, 159/13/1.3, 178/54/1.0, 196/21/1.6,
                        30/44/0.6, 74/57/1.2, 125/48/0.9, 168/6/0.8, 191/59/0.7,
                        52/18/1.0, 100/60/0.6, 148/20/1.1, 66/49/1.4, 205/45/0.8}
    \fill[Sky, opacity=0.30] ($(current page.north west)+(\x mm,-\y mm)$) circle (\r mm);
\end{tikzpicture}

\vspace*{4mm}
{\sffamily\fontsize{7.6}{9}\selectfont\color{Sky}%
 AUTOMATED DAILY LITERATURE DIGEST\;\textbullet\;PhD RESEARCH WATCH\par}
\vspace{2.2mm}
{\sffamily\bfseries\fontsize{27}{29}\selectfont\color{white}%
 Particulate Matter\par}
{\sffamily\fontsize{27}{29}\selectfont\color{Sky}%
 Monitoring \& Health Effects\par}
\vspace{5.0mm}
{\sffamily\fontsize{8.4}{10}\selectfont\color{white}%
 Issue 2026-08-30\;\;\textcolor{Amber}{\textbar}\;\;Entry window 30 August 2026 (single day)\;\;%
 \textcolor{Amber}{\textbar}\;\;14 records in scope, of which \textbf{3 are dated backfill}\par}
\vspace{18mm}

% -------------------------------------------------------------- metric strip
\noindent
\begin{minipage}[t]{0.190\linewidth}\metric{14}{RECORDS IN SCOPE\\(17 SCREENED OUT)}{Deep}\end{minipage}\hfill
\begin{minipage}[t]{0.190\linewidth}\metric{7}{SUBTOPIC CLUSTERS\\OF TEN OCCUPIED}{Teal}\end{minipage}\hfill
\begin{minipage}[t]{0.190\linewidth}\metric{4}{SENSING /\\INSTRUMENTATION}{Amber}\end{minipage}\hfill
\begin{minipage}[t]{0.190\linewidth}\metric{4}{EFFECT ESTIMATES\\(ONE SOURCE)}{Clay}\end{minipage}\hfill
\begin{minipage}[t]{0.190\linewidth}\metric{7}{RECORDS WITH NO\\HEALTH ENDPOINT}{Sky}\end{minipage}

\vspace{6mm}

% -------------------------------------------------------------- signal of day
\section*{\sffamily\bfseries\color{Deep}Signal of the day}
\vspace{-1.5mm}{\color{Amber}\rule{\linewidth}{1.1pt}}\vspace{2.5mm}

\begin{keybox}
{\sffamily\bfseries\small\color{Deep}A city-scale low-cost sensor network is used as the
\emph{instrument of record} for a mesoscale meteorology question --- and the calibration
literature published alongside it says why that is still risky.}\\[1.4mm]
\footnotesize
Chen et al.\ (\textit{Environ Sci Technol}) composite \textbf{42 lake-breeze events} over
Chicago in 2022 using one year of the Microsoft Eclipse low-cost sensor network plus
meteorological stations. Weekday lake breezes \emph{reduce} O$_3$ by up to
$\sim$6\,ppb along the northeast shoreline --- the opposite of the sea-breeze ozone
exceedances reported for most coastal regions --- while \emph{raising} PM$_{2.5}$ by up to
$\sim$3\,$\mu$g\,m$^{-3}$ through accumulation in a shallower boundary layer and enhanced
semivolatile partitioning at lower temperature. A 3\,$\mu$g\,m$^{-3}$ signal is precisely
the magnitude that the three calibration papers carried alongside it show is
\emph{recoverable but not free}: Blaga et al.\ report Pearson $r$ moving from
$-0.28$ to $0.89$ in the humid season after a two-step temperature/RH correction, and
from 0.06 to 0.65 in the dry season --- and the uncorrected Bucharest network flagged
exceedances at 8 of 44 sites where the corrected one flagged all 44. The methodological
reading for distributed PM sensing is blunt: an LCS network is now good enough to resolve
a few-$\mu$g\,m$^{-3}$ meteorological forcing, but only after a seasonally stratified,
site-specific correction, and Senarathna et al.\ show a wet-season model applied to
dry-season data carries \textbf{26.57\% MAPE} and does not transfer 115\,km between
Colombo and Kandy.
\end{keybox}

\vspace{3mm}

% -------------------------------------------------------------- exec brief
\section*{\sffamily\bfseries\color{Deep}What changed today}
\vspace{-1.5mm}{\color{Amber}\rule{\linewidth}{1.1pt}}\vspace{2.5mm}

\footnotesize
\begin{itemize}
\item \textbf{Calibration transferability is the binding constraint, stated three ways.}
Blaga (Bucharest, 44 Clarity Node-S), Senarathna (Sri Lanka, TSI BlueSky, two climatic
zones) and Sengupta (Oxford OxAria, LSTM) independently report the same failure mode:
a correction fitted at one site and season degrades elsewhere. Sengupta's contribution is
that a sequence model reaches \textbf{9.1\% expanded uncertainty for PM$_{2.5}$} against
the Equivalence Spreadsheet Tool --- inside the EU indicative-monitoring threshold --- by
learning time-lagged environmental dependence a Random Forest treating observations
independently cannot.

\item \textbf{The regulatory consequence of \emph{not} calibrating is now quantified.}
Bucharest's manufacturer pre-calibration understated PM$_{2.5}$ enough that only 18\% of
sensors appeared to breach the EU 25\,$\mu$g\,m$^{-3}$ daily limit; after correction every
site did, averaging $>$60 exceedance days against 35 permitted. This is the strongest
argument in the day's harvest for treating uncalibrated LCS output as unfit for
compliance framing.

\item \textbf{An operational null worth keeping.} Zhang \& Yu's Wuhan time-series
(906 days, 391{,}344 ED visits) recovers textbook associations --- PM$_{2.5}$ AQI
RR \textbf{1.028 (1.015--1.042)} per IQR at lag 0--2, PM$_{10}$ 1.030, NO$_2$ 1.049 ---
and then shows that \emph{adding those variables did not improve} day-ahead demand
forecasting (MAE 33.2 with meteorology, unchanged with pollution). Statistical detectability
and operational value are dissociable, and papers rarely test the second.

\item \textbf{Detection method, not environment, dominated airborne-microplastic
heterogeneity.} Rahmani et al.\ pool 27 studies across 16 countries at
$I^2>99\%$ and find the largest moderator is \emph{instrument} --- FTIR consistently
reporting higher counts than microscopy or Raman. Their indoor/outdoor contrast
(1{,}380 vs.\ 102\,MP\,m$^{-3}$) should not be read as an exposure ratio until sampling
and detection protocols are harmonised, and they say so.

\item \textbf{Two records ship on metadata alone.} The \textit{Aerosol Air Qual Res}
oil-mist ventilation study and the \textit{Environ Int} solid-fuel-PM nitrogen-organics
paper deposited no abstract to Crossref and are absent from Europe PMC, PubMed and
Semantic Scholar. They are listed, not summarised, and carry no extracted numbers.

\item \textbf{Health-side yield was low and skewed to synthesis.} Of 7 PubMed records
entering on 30 August, one was off-topic (flotation process water); the remainder are two
reviews, one animal study, one epigenetics cohort where air pollution is 1 of 25 mixture
components, one composition campaign, and one cross-sectional performance study.
No cardiovascular, respiratory or neuro-epidemiological record entered the window at all.
\end{itemize}

\clearpage

% ==================================================================== ANALYTICS
\section*{\sffamily\bfseries\color{Deep}Corpus analytics}
\vspace{-1.5mm}{\color{Amber}\rule{\linewidth}{1.1pt}}\vspace{3mm}

\noindent\includegraphics[width=\linewidth]{\FIGDIR/f1_subtopics.png}
\figcap{Subtopic assignment of the 14 in-scope records. Seven of the ten clusters are
occupied; \textit{Cardiovascular \& metabolic}, \textit{Neuro / mental health} and
\textit{Respiratory \& allergic} are all empty. Checked against \texttt{metrics.csv}: this is
the fourth issue this month with an empty respiratory cluster (21, 23 and 28 August were the
others), but the \emph{simultaneous} absence of respiratory, cardiovascular and neuro work is
unique in August. \textit{Sensing, forecasting \& instrumentation} leads at 4, but three of those
four are dated backfill recovered through Consensus and Crossref rather than deposits
entering PubMed on 30 August --- do not read this as a one-day flux.}

\vspace{2mm}
\noindent\includegraphics[width=\linewidth]{\FIGDIR/f2a_architecture.png}
\figcap{Study architecture. Sensor co-location is the single most common design (3),
followed by measurement campaigns (2) and reviews/syntheses (3, counting the microplastics
meta-analysis, the ILC2 review and the Southeast Asia scoping review). Two records are
\textit{Metadata only}: the architecture is genuinely unknown, not unmapped.}

\vspace{2mm}
\noindent\includegraphics[width=\linewidth]{\FIGDIR/f2b_endpoint.png}
\figcap{Health endpoint by organ system. \textbf{Seven of fourteen records carry no human
health endpoint} --- four sensing/calibration records, two composition or exposure-assessment
studies, and the metadata-only oil-mist record. Of the seven that do, no two share an
endpoint: emergency-department
attendance, mortality, athletic performance, neuroinflammation, cord-blood methylation,
in-vitro toxicology and inflammatory biomarkers each appear once. That dispersion is why no
pooled health estimate is attempted below.}

\vspace{2mm}
\noindent\includegraphics[width=\linewidth]{\FIGDIR/f3_metric_geography.png}
\figcap{Left --- particle metric. PM$_{2.5}$ alone is the exposure in 5 records; a further
3 report composition or speciation and 3 use unspeciated PM or emissions. Ultrafine
particles are again absent. Right --- geographic provenance. China contributes 3; the
single South Asian record (Sri Lanka, a calibration study) and the single Southeast Asian
record (the attribution review) are both present. Africa is absent today, but this is a
one-day gap and not a trend: the 29 August issue carried two African records
(South Africa; a multi-country sub-Saharan analysis).}

\vspace{2mm}
\noindent\includegraphics[width=\linewidth]{\FIGDIR/f6_lifecourse.png}
\figcap{Life-course windows addressed. Counts are non-exclusive and only five records
carry an identifiable window at all. The in-utero cell is the Barwon Infant Study cord-blood
analysis; the working-age cell holds the marathon cohort's 18--59 strata, the machining
workshop and the solid-fuel households; the 65+ cell is the marathon cohort's oldest
stratum, in which the NO$_2$ effect was smallest (1.32 vs.\ 2.17\,min per
$\mu$g\,m$^{-3}$ in 18--29-year-olds). Infancy/early childhood and adolescence are empty.}

\vspace{2mm}
\noindent\includegraphics[width=\linewidth]{\FIGDIR/f4_heatmap.png}
\figcap{Subtopic $\times$ study-architecture cross-tabulation. The structural point is the
same one this digest has recorded repeatedly: the \textit{Sensing} row is entirely
co-location and campaign work with no health-endpoint linkage, and the health rows contain
no instrumentation. Nothing in today's harvest bridges the two.}

\vspace{2mm}
\noindent\includegraphics[width=\linewidth]{\FIGDIR/f5_forest.png}
\figcap{All four ratio-metric estimates with confidence intervals in the corpus.
\textbf{They come from a single unrefereed preprint} (Zhang \& Yu, Research Square) and are
mutually adjusted co-pollutant estimates from one model, so the plot shows within-study
precision and direction only --- it is not a synthesis and carries no heterogeneity
statistic. The O$_3$ estimate is protective (0.980, 0.965--0.996), which in a Wuhan
time-series is most plausibly seasonal confounding rather than a causal effect.
Donaldson et al.'s marathon estimates are absolute time differences in minutes and are
deliberately excluded: on a log-ratio axis they would misread as null.}

\clearpage

% ==================================================================== DIGEST
\section*{\sffamily\bfseries\color{Deep}Paper-level digest}
\vspace{-1.5mm}{\color{Amber}\rule{\linewidth}{1.1pt}}\vspace{1mm}

{\fontsize{7.4}{9}\selectfont\color{Slate}Ordered by cluster. Each entry gives the
methodological core and the finding that matters, not an abstract paraphrase. Records
marked \textbf{[backfill]} were recovered on 30 August but were published earlier in 2026;
they are dated to this issue and must not be read as same-day deposits.\par}

\band{Sensing, forecasting \& instrumentation\hfill 4 records}{Amber}

\paperentry{Chen et al.\ 2026k}
{Lake-breeze effects on urban air quality resolved by a city-wide low-cost sensor network}
{Environ Sci Technol}
{One year (2022) of the Microsoft Eclipse LCS network plus meteorological stations over
Chicago, composited into \textbf{42 lake-breeze events} and 51 weak events confined within a
few kilometres of the shore. Weekday events reduce O$_3$ by up to $\sim$6\,ppb along the
northeast shoreline --- attributed to NO titration in a shallow boundary layer, inflow of
cleaner lake air, and suppressed photochemistry at lower temperature --- while
\emph{raising} PM$_{2.5}$ by up to $\sim$3\,$\mu$g\,m$^{-3}$ via accumulation of local
emissions and semivolatile partitioning. The limitation that matters: event-composited
means over an LCS network inherit that network's calibration bias, and no
reference-monitor validation of the PM$_{2.5}$ increment is reported, so the
3\,$\mu$g\,m$^{-3}$ figure is a network-relative anomaly rather than a traceable
concentration. Directly relevant precedent for using a dense LCS deployment as the primary
instrument in a mesoscale study.}
{10.1021/acs.est.6c09290}

\paperentry{Blaga et al.\ 2026 \textnormal{[backfill]}}
{City-scale calibration of a 44-node PM$_{2.5}$ network for regulatory-compliant assessment}
{PLOS Clim}
{44 Clarity Node-S sensors across Bucharest against National Environmental Protection Agency
reference stations; nine co-located spatial clusters, one used to fit and eight held out.
Manufacturer pre-calibration \textbf{systematically underestimated} PM$_{2.5}$, worst in
summer when secondary aerosol formation dominates. A two-step seasonal model on temperature
and relative humidity alone moved Pearson $r$ from \textbf{0.06 to 0.65} (dry season) and
from \textbf{$-$0.28 to 0.89} (humid season). Applied network-wide, calibrated sensors
average \textbf{$>$60 days/yr} above the EU 25\,$\mu$g\,m$^{-3}$ daily limit against 35
permitted; pre-calibration, only 8 of 44 (18\%) exceeded. Weakness: calibration degrades in
hot dry low-concentration conditions and inter-sensor variability in calibration performance
remains large, so a single network-wide model is not defensible. The 1:8 train/validate
split is the design detail worth copying.}
{10.1371/journal.pclm.0000780}

\paperentry{Senarathna et al.\ 2026 \textnormal{[backfill]}}
{Seasonal and spatial non-transferability of LCS calibration across two tropical climatic zones}
{Environ Monit Assess}
{TSI BlueSky sensors in a harmonised network across Colombo (wet-zone) and Kandy, compared
under linear regression, multiple linear regression and an echo-state network. Harmonised
sensors agreed within \textbf{10\%} of one another. Performance improves monotonically with
averaging window: at 24\,h, wet season $R^2=0.60$, MAE 0.93\,$\mu$g\,m$^{-3}$, MAPE 7.81\%;
dry season $R^2=0.85$, MAE 1.98, MAPE 9.17\%. The two transfer results are the point ---
applying a wet-season model to dry-season data yields \textbf{26.57\% MAPE}, and a
Colombo-fitted model degrades in Kandy, while a Kandy-specific model with temperature and RH
reaches $R^2=0.84$--$0.92$. Limitation: no independent reference co-location is reported for
the second zone at the same duration, so the spatial-transfer claim rests partly on
inter-sensor agreement rather than reference truth.}
{10.1007/s10661-026-15623-4}

\paperentry{Sengupta et al.\ 2026 \textnormal{[backfill]}}
{LSTM calibration of low-cost PM and NO$_2$ sensors against a regulatory equivalence test}
{arXiv (preprint)}
{Sequence model trained on co-located OxAria reference data in Oxford for PM$_{2.5}$,
PM$_{10}$ and NO$_2$, benchmarked against a Random Forest baseline. The framing argument is
that RF treats each hourly observation independently and therefore cannot represent sensor
lag or delayed environmental response; the LSTM with time-lagged features, harmonic time
encodings and interaction terms beats it on $R^2$ across train, validation and test.
The result that carries weight is the external check: validation against Equivalence
Spreadsheet Tool 3.1 gives expanded uncertainties of \textbf{9.1\% (PM$_{2.5}$),
12.42\% (PM$_{10}$) and 22.11\% (NO$_2$)} --- i.e.\ PM meets, and NO$_2$ misses, the
indicative-monitoring bar. Unrefereed, single network, and no drift analysis over a
multi-year horizon, which is where sequence models are most likely to fail.}
{10.48550/arXiv.2604.21527}

\band{Exposure assessment \& modelling\hfill 2 records}{Amber}

\paperentry{Rahmani et al.\ 2026}
{Airborne microplastics indoors and outdoors: pooled concentrations and the instrument problem}
{Rev Environ Health}
{PRISMA systematic review to October 2023: 27 eligible studies across 16 countries, 21
pooled under a random-effects model. Indoor concentrations exceed outdoor --- moderately
polluted indoor environments 1{,}380.28 vs.\ 102.45\,MP\,m$^{-3}$ outdoors --- with indoor
healthcare settings and outdoor urban areas the most contaminated sub-environments.
The authors present the pooled figures as a descriptive summary rather than an estimate,
correctly: heterogeneity is $\mathbf{I^2>99\%}$ and the dominant moderator is
\textbf{detection instrument}, with FTIR reporting consistently higher counts than
microscopy or Raman spectroscopy; sampling duration also matters. Sensitivity analysis
excludes single-outlier dominance. For instrumentation work this is a clean external
demonstration that method-driven variance can exceed environmental variance, which is
exactly the risk in heterogeneous LCS fleets.}
{10.1515/reveh-2025-0059}

\paperentry{Ren et al.\ 2026}
{Vertical contrast in PM$_{2.5}$ composition and secondary formation, Shiyan, central China}
{iScience}
{One year across three urban sites and one mountain site. Annual mountain PM$_{2.5}$ is
\textbf{70\% of urban} levels, but sulfate--nitrate--ammonium is a \emph{larger} fraction
there (49.3\%). SNA is vertically near-homogeneous in spring, summer and autumn
(coefficient of divergence $<$0.2) and stratifies only in winter, when most components are
urban-enriched. Higher nitrogen oxidation ratio (0.27) and SOC/OC (0.65) at altitude
indicate a more aged aerosol. Limitation: a single mountain site cannot separate altitude
from the site's own upwind exposure, and no vertical profiling instrument is used --- the
"vertical distribution" claim rests on a two-point contrast. Relevant to siting logic for
distributed networks: composition, not mass, is what varies with height here, and a
mass-only node would report the mountain site as simply cleaner.}
{10.1016/j.isci.2026.117090}

\band{Occupational \& indoor\hfill 2 records}{Clay}

\paperentry{Gu et al.\ 2026b}
{Nitrogen-containing organics in household solid-fuel combustion PM vs.\ urinary inflammatory
and amino-acid profiles}
{Environ Int}
{Listed on Crossref metadata only. The deposit carries no abstract, and the record is absent
from Europe PMC, PubMed and Semantic Scholar as of this run, so \textbf{no finding, effect
size, cohort size or location is asserted here}. The title indicates an exploratory personal
exposure design pairing speciated N-containing organic compounds in solid-fuel combustion PM
with targeted urinary inflammatory and amino-acid metabolites --- a personal-monitoring plus
targeted-metabolomics pairing that is directly on-axis for this watch. Queued for the
abstract retry list; it will be summarised, not re-listed, once text is available.}
{10.1016/j.envint.2026.110491}

\paperentry{Li et al.\ 2026f}
{Non-uniform oil-mist distribution and local ventilation in a machining workshop}
{Aerosol Air Qual Res}
{Listed on Crossref metadata only --- no abstract deposited, and no record in Europe PMC,
PubMed or Semantic Scholar. \textbf{No quantitative claim is made here.} The title indicates
spatial characterisation of oil-mist aerosol in a machining workshop with a local
ventilation intervention, which is the occupational analogue of the indoor-siting problem:
where a workshop's aerosol field is strongly non-uniform, a single fixed monitor is a poor
estimator of personal exposure regardless of its accuracy. Queued for the abstract retry
list.}
{10.1007/s44408-026-00161-y}

\band{Mechanistic toxicology\hfill 2 records}{Violet}

\paperentry{Jerome et al.\ 2026}
{World Trade Center dust: neuroinflammation and myelin response at 3 and 12 months}
{J Immunotoxicol}
{Twelve-week-old male spontaneously hypertensive rats given intratracheal WTC dust at
\textbf{2.43 $\pm$ 0.78\,mg/day on two consecutive days}, then euthanised at 3 and 12 months.
At 3 months, increased Black-Gold II myelin staining without change in myelin basic protein
or CNPase; at 12 months, microglial activation (Iba1) and elevated brain
\textbf{IL-1$\beta$}. The authors read the discrepant myelin markers as a compensatory
response preceding degeneration. Three limitations bound the inference: a two-day bolus
instillation is not a chronic inhalation exposure and bypasses the upper airway; the SHR
strain is hypertensive by design, so the neuroinflammatory phenotype is confounded with a
vascular one; and males only. Useful as a mechanistic anchor for the First-Responder
cognitive literature, not as a dose--response.}
{10.1080/1547691X.2026.2722654}

\paperentry{Luo et al.\ 2026b}
{Environmental pollutants and innate immunity via group 2 innate lymphoid cells}
{Clin Rev Allergy Immunol}
{Narrative review covering O$_3$, PM, diesel exhaust particles and cigarette smoke alongside
nanomaterials and microplastics, organised around ILC2 biology. The proposed common route is
epithelial barrier compromise followed by alarmin release --- IL-33, IL-25, TSLP --- with
ILC2s as the sensor coupling environmental detection to type-2 inflammation, plus
non-canonical regulation of ILC2 plasticity and functional adaptation. No new data, no
quantitative synthesis, and PM is not separated from the other exposures, so nothing here
constrains a PM-specific dose or particle-property relationship. Carried as a citation
anchor for the epithelial-alarmin axis in respiratory PM work; tier C.}
{10.1007/s12016-026-09193-z}

\band{Reproductive \& developmental\hfill 1 record}{Sky}

\paperentry{Tanner et al.\ 2026}
{Prenatal exposure mixture and aromatase brain-promoter methylation in cord blood}
{Environ Epigenet}
{Barwon Infant Study birth cohort, \textbf{n = 906}, testing 25 prenatal exposures across
five classes --- manufactured chemicals, air pollution, pharmacological, nutritional and
sunlight-related --- against methylation of the brain-specific \textit{CYP19A1} P1.f promoter
using weighted quantile sum regression. The mixture index was associated with higher P1.f
methylation, \textbf{adjusted mean difference 0.71 (95\% CI 0.11--1.32), $p=0.021$},
interpreted as reduced brain aromatase activity. Bisphenols, reduced sunlight, household
mould, phthalates, low folate and air pollution carried the largest weights. The
inference-threatening limitation is structural: WQS assigns weights within a single index,
so air pollution's individual contribution is neither estimated nor separable, and cord
blood is a proxy for a brain-specific promoter. Tier C for this watch --- PM is one component
of a 25-exposure mixture.}
{10.1093/eep/dvag027}

\band{Burden, policy \& mitigation\hfill 1 record}{Coral}

\paperentry{Simion et al.\ 2026}
{Climate-related health attribution in Southeast Asia: what is and is not being studied}
{Lancet Reg Health West Pac}
{Scoping review, \textbf{5{,}294 records screened, 77 included}. Air pollution is the subject
of \textbf{70\% of studies}, linked to up to 31\% of all-cause mortality and an estimated
\textbf{296{,}000 annual deaths from wildfire smoke during El Ni\~no events}; heat is second,
associated with up to 8.61\% of mortality in Bangkok. Floods, droughts, extreme rainfall,
storms and sea-level rise are near-absent. The paper's own framing --- that the research
portfolio is misaligned with the regional hazard profile --- is the transferable finding; the
attributable fractions themselves are heterogeneously derived across 77 studies with no
standardised attribution framework, which the authors name as the central methodological gap.
For a monitoring programme the operative point is that the region's dominant PM source is
episodic biomass burning, which mass-only fixed networks under-sample in both space and time.}
{10.1016/j.lanwpc.2026.101952}

\band{Other clinical endpoints\hfill 2 records}{Slate}

\paperentry{Donaldson et al.\ 2026}
{Race-day air pollution and marathon finishing time across 2.76 million finishes}
{Sports Med Open}
{\textbf{2{,}756{,}553} net finishing times from six World Marathon Majors (Berlin, Boston,
Chicago, London, New York, Tokyo), 75 marathon-year events, 2010--2024, with race-day NO$_2$
and PM$_{2.5}$ from CAMS global reanalysis and linear mixed models adjusting heat index, wind
and temporal trend. Each 1\,$\mu$g\,m$^{-3}$ NO$_2$ was associated with
\textbf{+1.43\,min (95\% CI 0.66--2.21) in men and +1.56\,min (0.63--2.48) in women}, with a
steep gradient by ability --- recreational 2.12 vs.\ elite 0.34\,min per $\mu$g\,m$^{-3}$,
persisting as a percentage of mean time (0.67\% vs.\ 0.18\%) --- and by age
(2.17\,min in 18--29-year-olds vs.\ 1.32 in $\geq$60). \textbf{PM$_{2.5}$ was null}
($-$0.13\,min, $-$0.67 to 0.40 in men). The exposure is reanalysis at city scale, not
course-level or personal, and NO$_2$ at that resolution is a traffic-proximity proxy; the
NO$_2$/PM$_{2.5}$ divergence is therefore at least partly a measurement-error contrast
rather than a toxicological one. Still the largest exercise-performance exposure dataset
assembled.}
{10.1186/s40798-026-01099-6}

\paperentry{Zhang \& Yu 2026 (preprint)}
{Air pollution is detectable in ED demand and useless for forecasting it}
{Res Sq (preprint)}
{Single-centre daily time-series at Tongji Hospital Guanggu Branch, Wuhan, 1 Jan 2024 to
26 Jun 2026: \textbf{906 dates, 391{,}344 visits}, mean 431.9/day (SD 59.2), falling to
319.3/day over Spring Festival. Quasi-Poisson models give, per IQR of lag 0--2
pollutant-specific AQI, RR \textbf{1.028 (1.015--1.042)} for PM$_{2.5}$,
1.030 (1.017--1.043) for PM$_{10}$, 1.049 (1.034--1.064) for NO$_2$ and an inverse
0.980 (0.965--0.996) for O$_3$; each 1\,$^\circ$C is +0.85\% visits. The contribution is the
temporally-ordered validation: a history-plus-calendar model achieves MAE 33.6 visits/day,
meteorology reduces it to 33.2, and \textbf{adding air-pollution AQI does not improve it at
all}. Single centre, city-level AQI rather than catchment exposure, unrefereed. The inverse
O$_3$ term is almost certainly seasonal confounding. Read as a warning that a statistically
robust exposure--response can carry no incremental predictive information.}
{10.21203/rs.3.rs-10496137/v1}

\clearpage

% ==================================================================== TABLE
\section*{\sffamily\bfseries\color{Deep}Corpus register}
\vspace{-1.5mm}{\color{Amber}\rule{\linewidth}{1.1pt}}\vspace{2.5mm}

{\fontsize{6.7}{8.4}\selectfont
\setlength{\tabcolsep}{2.4pt}
\begin{longtable}{@{}>{\RaggedRight}p{25mm}>{\RaggedRight}p{41mm}>{\RaggedRight}p{27mm}%
>{\RaggedRight}p{23mm}>{\RaggedRight}p{18mm}>{\RaggedRight}p{15mm}@{}}
\toprule
\textsf{\textbf{First author}} & \textsf{\textbf{Focus}} & \textsf{\textbf{Journal}} &
\textsf{\textbf{Design}} & \textsf{\textbf{Metric}} & \textsf{\textbf{Region}}\\
\midrule
\endfirsthead
\toprule
\textsf{\textbf{First author}} & \textsf{\textbf{Focus}} & \textsf{\textbf{Journal}} &
\textsf{\textbf{Design}} & \textsf{\textbf{Metric}} & \textsf{\textbf{Region}}\\
\midrule
\endhead
\multicolumn{6}{@{}l@{}}{\cellcolor{Mist}\textsf{\textbf{\textcolor{Deep}{Neuro / mental health}}}}\\[0.6mm]
\href{https://doi.org/10.1001/jamanetworkopen.2026.25015}{Dales} & Wildfire-sourced PM$_{2.5}$ and ED visits for migraine and other primary headache & \textit{JAMA Netw Open} & Case-crossover & PM$_{2.5}$ (wildfire) & Canada\\
\href{https://doi.org/10.1093/ageing/afag215}{Glover} & Umbrella meta-review of urban environmental exposures and cognitive health & \textit{Age Ageing} & Umbrella review / meta-analysis & PM$_{2.5}$ & Global\\
\href{https://doi.org/10.1016/j.nsa.2026.107019}{Goossens} & In utero air pollution exposure and pre-/postnatal brain development & \textit{Neurosci Appl} & Narrative review & PM$_{2.5}$ / PM$_{10}$ / PM$_{\mathrm{coarse}}$ & Global\\
\href{https://doi.org/10.1016/j.isci.2026.116617}{Ma} & PM$_{2.5}$-induced neurotoxicity and depression: exposure risks and mechanisms & \textit{iScience} & Narrative review & PM$_{2.5}$ + components & Global\\
\href{https://doi.org/10.3760/cma.j.cn112150-20250804-00752}{Pan} & Maternal PM exposure and offspring neurodevelopment: effects and mechanisms & \textit{Zhonghua Yu Fang Yi Xue Za Zhi} & Narrative review & PM (bulk) & China\\
\addlinespace[0.9mm]
\multicolumn{6}{@{}l@{}}{\cellcolor{Mist}\textsf{\textbf{\textcolor{Teal}{Cardiovascular \& metabolic}}}}\\[0.6mm]
\href{https://doi.org/10.3760/cma.j.cn112338-20251120-00843}{Huang} & PM$_{2.5}$ and dyslipidaemia in rural Chinese children; dietary effect modification & \textit{Zhonghua Liu Xing Bing Xue Za Zhi} & Cross-sectional (multistage) & PM$_{2.5}$ & China\\
\href{https://doi.org/10.32394/pe/218966}{Kazmierska} & Air pollution and cardiovascular disease: review of current data & \textit{Przegl Epidemiol} & Narrative review & PM$_{2.5}$ / PM$_{10}$ & Global\\
\href{https://doi.org/10.1007/s00125-026-06814-2}{Xi} & Air pollution, built environment and temperature vs. prediabetes phenotypes & \textit{Diabetologia} & Prospective cohort & PM$_{2.5}$ & Germany\\
\href{https://doi.org/10.1161/JAHA.125.046917}{Yang} & Short-term PM$_{2.5}$ and hospitalisation/ED use in atrial fibrillation & \textit{J Am Heart Assoc} & Case-crossover & PM$_{2.5}$ (wildfire-influenced) & USA\\
\addlinespace[0.9mm]
\multicolumn{6}{@{}l@{}}{\cellcolor{Mist}\textsf{\textbf{\textcolor{Sky}{Reproductive \& developmental}}}}\\[0.6mm]
\href{https://doi.org/10.1016/j.jes.2025.10.043}{Chen} & Prenatal air pollution, fetal growth trajectory phenotypes and child neurodevelopment & \textit{J Environ Sci (China)} & Prospective cohort & PM$_{2.5}$ / PM$_{10}$ & China\\
\href{https://doi.org/10.1016/j.ecoenv.2026.120566}{Huang} & Joint prenatal bisphenol + air-pollutant mixture exposure and preschool IQ & \textit{Ecotoxicol Environ Saf} & Prospective cohort (mixtures) & PM$_{2.5}$ / PM$_{10}$ & China\\
\href{https://doi.org/10.1007/s40572-026-00552-8}{Peterson} & Ambient air pollution, PAHs and pubertal development: critical review & \textit{Curr Environ Health Rep} & Critical review & PM$_{2.5}$ components & Global\\
\href{https://doi.org/10.1016/j.envres.2026.125313}{Xie} & Long-term PM$_{2.5}$ and timing of pubertal development in boys and girls & \textit{Environ Res} & Prospective cohort & PM$_{2.5}$ & China\\
\addlinespace[0.9mm]
\multicolumn{6}{@{}l@{}}{\cellcolor{Mist}\textsf{\textbf{\textcolor{Sage}{Respiratory \& allergic}}}}\\[0.6mm]
\href{https://doi.org/10.1007/s00484-026-03289-z}{Bayar Muluk} & Climate change and otolaryngology: airway disease impacts and emission reduction & \textit{Int J Biometeorol} & Narrative review & PM$_{2.5}$ & Global\\
\href{https://doi.org/10.4103/mgr.MEDGASRES-D-25-00164}{Liao} & Short-term air pollution, respiratory symptoms and lung function in asthma & \textit{Med Gas Res} & Panel study & PM$_{2.5}$ / PM$_{10}$ & China\\
\href{https://doi.org/10.3390/epidemiologia7040097}{Sanchez de Prada} & Ambient PM and asthma-coded admissions in children 0-3 y, Spain & \textit{Epidemiologia} & Ecological panel & PM$_{2.5}$ / PM$_{10}$ & Spain\\
\href{https://doi.org/10.1016/j.jhazmat.2026.143076}{Yang} & Allergenic vs. toxic risk of airborne seed hairs compared with ambient PM$_{2.5}$ & \textit{J Hazard Mater} & Field measurement & PM$_{2.5}$ + bioaerosol & China\\
\addlinespace[0.9mm]
\multicolumn{6}{@{}l@{}}{\cellcolor{Mist}\textsf{\textbf{\textcolor{Violet}{Mechanistic toxicology}}}}\\[0.6mm]
\href{https://doi.org/10.1016/j.jes.2025.12.074}{Ding} & PM$_{2.5}$ promotes vascular calcification via miR-27b-5p/HES1 signalling & \textit{J Environ Sci (China)} & In vitro + organoid + murine & PM$_{2.5}$ & China\\
\href{https://doi.org/10.1007/s00360-026-01698-5}{Maraschi} & Osmoregulatory disruption and metal bioaccumulation from settleable atmospheric PM & \textit{J Comp Physiol B} & Controlled exposure (ecotox) & Settleable APM & Brazil\\
\href{https://doi.org/10.1016/j.jes.2025.11.018}{Wu} & Macrophage NLRP3 inflammasome in real-ambient PM-induced hepatic glucose dysregulation & \textit{J Environ Sci (China)} & Real-ambient murine + co-culture & Real-ambient PM & China\\
\addlinespace[0.9mm]
\multicolumn{6}{@{}l@{}}{\cellcolor{Mist}\textsf{\textbf{\textcolor{Clay}{Occupational \& indoor}}}}\\[0.6mm]
\href{https://doi.org/10.3389/fmedt.2026.1796828}{Chen} & Nested suction instrument to cut intraoperative surgical-smoke PM$_{2.5}$ & \textit{Front Med Technol} & Bench / device experiment & PM$_{2.5}$ (surgical smoke) & China\\
\href{https://doi.org/10.1093/annweh/wxag055}{Cole} & Determinants of respirable crystalline silica exposure in Australian tunnelling & \textit{Ann Work Expo Health} & Observational exposure study & Respirable dust / RCS & Australia\\
\href{https://doi.org/10.1016/j.jhazmat.2026.143081}{Ding} & Stage-resolved hair Zn-isotope dynamics as a firefighting combustion-exposure tracer & \textit{J Hazard Mater} & Repeated-measures biomarker & Combustion-derived PM & China\\
\href{https://doi.org/10.3390/jox16040119}{Ramsden} & Xenobiotic hazards in aircraft cabin air & \textit{J Xenobiot} & Review + risk economics & Ultrafine particles & Global\\
\addlinespace[0.9mm]
\multicolumn{6}{@{}l@{}}{\cellcolor{Mist}\textsf{\textbf{\textcolor{Amber}{Sensing, forecasting \& instrumentation}}}}\\[0.6mm]
\href{https://doi.org/10.1016/j.jenvman.2026.130515}{Shi} & Adaptive spatiotemporal graph transformer for multi-site multi-step PM$_{2.5}$ forecasting & \textit{J Environ Manage} & Deep-learning model & PM$_{2.5}$ & China / India\\
\href{https://doi.org/10.3390/bios16070380}{Tsuchiya} & Seasonal and indoor-outdoor characteristics of fluorescent aerosol particles in offices & \textit{Biosensors} & Real-time sensor campaign & Bioaerosol / coarse & Japan\\
\addlinespace[0.9mm]
\multicolumn{6}{@{}l@{}}{\cellcolor{Mist}\textsf{\textbf{\textcolor{Amber}{Exposure assessment \& modelling}}}}\\[0.6mm]
\href{https://doi.org/10.1029/2025GH001621}{Solorzano-Araque} & Multiscale exposure modelling: LOTOS-EUROS + mobility surveys + particle toxicity & \textit{GeoHealth} & CTM + mobility integration & PM$_{2.5}$ / PM$_{10}$ & Colombia\\
\addlinespace[0.9mm]
\multicolumn{6}{@{}l@{}}{\cellcolor{Mist}\textsf{\textbf{\textcolor{Coral}{Burden, policy \& mitigation}}}}\\[0.6mm]
\href{https://doi.org/10.1002/hsr2.72886}{Karamzad} & Ischaemic heart disease burden in MENA 1990-2021; PM attributable fraction 28.2\% & \textit{Health Sci Rep} & GBD secondary analysis & Ambient PM & MENA\\
\href{https://doi.org/10.1371/journal.pone.0347989}{Rees} & HEPA filtration to reduce respiratory infection in care homes & \textit{PLoS One} & RCT process evaluation & Filtered PM & UK\\
\href{https://doi.org/10.1016/j.jes.2026.01.009}{Shan} & Policy-driven emission reductions and post-pandemic rebound, high-resolution inventories & \textit{J Environ Sci (China)} & Emission inventory & PM emissions & China\\
\addlinespace[0.9mm]
\multicolumn{6}{@{}l@{}}{\cellcolor{Mist}\textsf{\textbf{\textcolor{Slate}{Other clinical endpoints}}}}\\[0.6mm]
\href{https://doi.org/10.3390/dj14070440}{Lupita} & PM$_{2.5}$, preventive-care capacity and caries/edentulism burden across the EU & \textit{Dent J} & Ecological panel & PM$_{2.5}$ & EU-27\\
\href{https://doi.org/10.1016/j.jdin.2026.06.008}{Park} & Long-term PM exposure and risk of androgenetic alopecia & \textit{JAAD Int} & Nationwide cohort & PM$_{2.5}$ / PM$_{10}$ & South Korea\\
\href{https://doi.org/10.3171/2026.3.PEDS2615}{Sun} & Environmental exposure and area deprivation vs. paediatric hydrocephalus outcomes & \textit{J Neurosurg Pediatr} & Registry cohort + geospatial & PM$_{2.5}$ (satellite-derived) & USA\\
\addlinespace[0.9mm]
\bottomrule
\end{longtable}}


\vspace{4mm}
\begin{keybox}
{\sffamily\bfseries\footnotesize\color{Deep}Provenance and method}\\[1.2mm]
{\fontsize{7.2}{8.8}\selectfont
\textbf{Window: 30 August 2026 alone}, contiguous with the 29 August issue; no gap.
Sources queried on 31 August 2026. \textbf{PubMed} \texttt{[EDAT]} 2026/08/30 --- health
axis 7 via the connector against 4 from the local harvester (the connector is again the
better recall leg; the local set is a strict subset), instrumentation axis \textbf{0} on
both. \textbf{Europe PMC} \texttt{CREATION\_DATE} 2026-08-30 --- 6 hits, of which 2 were
already in the PubMed set and 2 were published in the 29 August issue, leaving 2 unique
(one review, one Research Square preprint). \textbf{Crossref by-ISSN} across 25 tracked
non-PubMed journals --- 12 deposits with a 30 August \texttt{created} date, of which 3 are
carried and 9 screened out as non-particulate (PFAS, PCB, soil vanadium, gas-phase ozone
isotopes, aquatic pharmacology, biodiversity accounting). \textbf{arXiv} relevance query
returned 60 entries, \textbf{none} published on or after 25 August. \textbf{OpenAlex}
remains dead --- HTTP 403 paid-plan block with intermittent 429 --- and is routed through
Consensus, which returned 10 hits on the calibration/drift axis: 7 resolved to DOIs already
in \texttt{seen.json}, \textbf{3 were new} and are carried as explicitly labelled backfill
(Blaga, Senarathna, Sengupta; publication dates February, June and April 2026).
\textbf{31 unique records considered; 7 dropped as already seen; 14 in scope; 17 logged to
\texttt{rejected.jsonl} with a reason.}
\emph{Two records are listed on metadata only} --- \texttt{10.1016/j.envint.2026.110491} and
\texttt{10.1007/s44408-026-00161-y} deposited no abstract to Crossref and are absent from
Europe PMC, PubMed and Semantic Scholar; they are not summarised and no numbers are
attributed to them.
Subtopic, design, particle-metric and geography labels are assigned from title, abstract,
keywords and MeSH terms; assignment is single-label by dominant contribution and is a
judgement, not a controlled vocabulary. First-author surnames are read from PubMed
\texttt{efetch} or the Crossref deposit, never inferred.
All 14 DOIs were checked against the PMID's own PubMed record and resolved at
\texttt{doi.org} before this issue was built: \textbf{0 failures, 0 advisory warnings}.
Content derived from PubMed and Crossref metadata; abstracts remain the copyright of their
respective publishers.\par}
\end{keybox}

\end{document}
